\input{../macros/preambulo.tex}

\begin{document}

% --- PORTADA ---
\begin{center}
    \Huge \textbf{Apunte Único: Algoritmos y Estructuras de Datos - Guía 1} \\
    \vspace{0.5cm}
    \Large Por alumnos de AED \\
    \normalsize Facultad de Ciencias Exactas y Naturales \\
    UBA \\
    \vspace{0.3cm}
    \small Última actualización: \today
\end{center}
\vspace{1cm}

% --- ÍNDICE ---
\hypertarget{indice}{}
\noindent\textbf{\Large The Council Will Decide Your Fate:} \\
\textit{(doble click en los ejercicios para saltar)}

\begin{multicols}{4}
\begin{itemize}
    \item \hyperlink{ej1}{Ejercicio 1}
    \item \hyperlink{ej2}{Ejercicio 2}
    \item \hyperlink{ej3}{Ejercicio 3}
    \item \hyperlink{ej4}{Ejercicio 4}
    \item \hyperlink{ej5}{Ejercicio 5}
    \item \hyperlink{ej6}{Ejercicio 6}
    \item \hyperlink{ej7}{Ejercicio 7}
    \item \hyperlink{ej8}{Ejercicio 8}
    \item \hyperlink{ej9}{Ejercicio 9}
    \item \hyperlink{ej10}{Ejercicio 10}
    \item \hyperlink{ej11}{Ejercicio 11}
    \item \hyperlink{ej12}{Ejercicio 12}
\end{itemize}
\end{multicols}

\vspace{0.5cm}


\begin{disclaimerbox}
\textbf{Disclaimer:} Recomendable leer lo siguiente, ayuda banda si no estas preparado para encarar la guía usando el repo. \\

Estudiar con resueltos puede ser un arma de doble filo. Si estás trabado, antes de saltar a la solución que hizo otra persona:
\begin{enumerate}
    \item \textbf{No mires la solución ni bien te trabás.} Te condicionás pavlovianamente a no pensar. Necesitás darle tiempo al cerebro.
    \item Intentá un ejercicio similar, pero más fácil. ¿No sale? Intentá uno aún más fácil.
    \item Fijate si tenés un ejercicio similar hecho en clase.
    \item Tomate 2 minutos para formular una pregunta real sobre lo que no entendés.
\end{enumerate}
Eh loco, relajá un toque... ¡va a salir todo bien!
\end{disclaimerbox}

\newpage

% --- COMIENZO DE LA GUÍA ---
\section*{1.1. Repaso de lógica proposicional}

% --- ejercicio_1.tex ---
\hypertarget{ej1}{}
\noindent \textbf{Ejercicio 1.} Determinar los valores de verdad de las siguientes proposiciones cuando el valor de verdad de $a$, $b$ y $c$ es verdadero y el de $x$ e $y$ es falso.

\begin{multicols}{2}
\begin{enumerate}[label=\alph*)]
    \item $(\neg x \vee b)$
    \item $((c \vee (y \wedge a)) \vee b)$
    \item $\neg(c \vee y)$
    \item $\neg(y \vee c)$
    \item $(\neg(c \vee y) \leftrightarrow (\neg c \wedge \neg y))$
    \item $((c \vee y) \wedge (a \vee b))$
    \item $(((c \vee y) \wedge (a \vee b)) \leftrightarrow (c \vee (y \wedge a) \vee b))$
    \item $(\neg c \wedge \neg y)$
\end{enumerate}
\end{multicols}

\vspace{0.2cm}
{\color{gray}\hrule}
\vspace{0.4cm}

\textbf{Resolución:}

\textbf{a)} Para dar una demostración del valor de verdad de esta proposición, primero identificamos 
los valores dados por el enunciado. Sabemos que el valor de $b$ es verdadero ($V$) y el de $x$ es falso ($F$).

Como ya tenemos sus valores fijos, basta con armar una pequeña tabla de valores paso a paso. Primero calculamos la negación de $x$ ($\neg x$) y luego resolvemos la disyunción lógica ($\vee$) con $b$:

\begin{center}
\begin{tabular}{|c|c|c|c|}
\hline
$x$ & $\neg x$ & $b$ & $(\neg x \vee b)$ \\
\hline
F & V & V & \textbf{V} \\
\hline
\end{tabular}
\end{center}

Y así comprobamos que la proposición $(\neg x \vee b)$ es \textbf{Verdadera}.

\vspace{0.4cm}
\textbf{b)} Para resolver esta proposición, primero reemplazamos los valores conocidos: $c=V$, $y=F$, $a=V$ y $b=V$. Al igual que en matemática, resolvemos de adentro hacia afuera, comenzando por el paréntesis más interno $(y \wedge a)$:

\begin{center}
\begin{tabular}{|c|c|c|c|c|c|c|}
\hline
$c$ & $y$ & $a$ & $b$ & $(y \wedge a)$ & $(c \vee (y \wedge a))$ & $((c \vee (y \wedge a)) \vee b)$ \\
\hline
V & F & V & V & F & V & \textbf{V} \\
\hline
\end{tabular}
\end{center}

Y así comprobamos que la proposición es \textbf{Verdadera}.

\vspace{0.4cm}
\textbf{c)} Reemplazamos $c=V$ e $y=F$. Primero evaluamos la disyunción dentro del paréntesis y luego aplicamos la negación al resultado final:

\begin{center}
\begin{tabular}{|c|c|c|c|}
\hline
$c$ & $y$ & $(c \vee y)$ & $\neg(c \vee y)$ \\
\hline
V & F & V & \textbf{F} \\ 
\hline
\end{tabular}
\end{center}

Por lo tanto, concluimos que la proposición es \textbf{Falsa}.

\vspace{0.4cm}
\textbf{d)} Como la disyunción lógica es conmutativa ($c \vee y \equiv y \vee c$), esta expresión es lógicamente equivalente a la del inciso anterior. Lo comprobamos armando la tabla con $y=F$ y $c=V$:

\begin{center}
\begin{tabular}{|c|c|c|c|}
\hline
$y$ & $c$ & $(y \vee c)$ & $\neg(y \vee c)$ \\
\hline
F & V & V & \textbf{F} \\
\hline
\end{tabular}
\end{center}

Y así comprobamos que, efectivamente, la proposición también es \textbf{Falsa}.

\vspace{0.4cm}
\textbf{e)} Aquí estamos frente a una de las Leyes de De Morgan, por lo que sabemos de antemano que debería resultar en una tautología (verdadera sin importar los valores). Lo demostramos reemplazando con nuestros valores dados $c=V$ e $y=F$:

\begin{center}
\begin{tabular}{|c|c|c|c|c|c|c|}
\hline
$c$ & $y$ & $\neg(c \vee y)$ & $\neg c$ & $\neg y$ & $(\neg c \wedge \neg y)$ & $(\neg(c \vee y) \leftrightarrow (\neg c \wedge \neg y))$ \\
\hline
V & F & F & F & V & F & \textbf{V} \\
\hline
\end{tabular}
\end{center}

Al tener el mismo valor de verdad de ambos lados ($F \leftrightarrow F$), comprobamos que la proposición es \textbf{Verdadera}.

\vspace{0.4cm}
\textbf{f)} Reemplazamos $c=V$, $y=F$, $a=V$ y $b=V$. Evaluamos de forma independiente las dos disyunciones de los paréntesis y luego las unimos con la conjunción principal:

\begin{center}
\begin{tabular}{|c|c|c|c|c|c|c|}
\hline
$c$ & $y$ & $a$ & $b$ & $(c \vee y)$ & $(a \vee b)$ & $((c \vee y) \wedge (a \vee b))$ \\
\hline
V & F & V & V & V & V & \textbf{V} \\
\hline
\end{tabular}
\end{center}

Como ambos lados de la conjunción son verdaderos, la proposición es \textbf{Verdadera}.

\vspace{0.4cm}
\textbf{g)} Si observamos con atención, el lado izquierdo de esta equivalencia es exactamente la proposición que evaluamos en el inciso \textbf{f)}, y el lado derecho es la que evaluamos en el inciso \textbf{b)}. Para ahorrar tiempo y no hacer una tabla gigante, podemos usar esos resultados directamente:

\begin{center}
\begin{tabular}{|c|c|c|} 
\hline
Lado Izquierdo (inciso f) & Lado Derecho (inciso b) & Equivalencia ($\leftrightarrow$) \\
\hline
V & V & \textbf{V} \\
\hline
\end{tabular}
\end{center}

Como ambos lados evalúan a verdadero ($V \leftrightarrow V$), comprobamos que la proposición final es \textbf{Verdadera}.

\vspace{0.4cm}
\textbf{h)} Reemplazamos $c=V$ e $y=F$. Evaluamos primero sus respectivas negaciones y finalmente aplicamos la conjunción:

\begin{center}
\begin{tabular}{|c|c|c|c|c|}
\hline
$c$ & $y$ & $\neg c$ & $\neg y$ & $(\neg c \wedge \neg y)$ \\
\hline
V & F & F & V & \textbf{F} \\
\hline
\end{tabular}
\end{center}

La conjunción requiere que ambos términos sean verdaderos para ser verdadera. Como $\neg c$ es falso, comprobamos que la proposición es \textbf{Falsa}.

\vspace{0.5cm}
\hrule
\vspace{0.5cm}
\hypertarget{ej2}{}
\noindent \textbf{Ejercicio 2.} Considere la siguiente oración: ``Si es mi cumpleaños o hay torta, entonces hay torta''.
\begin{itemize}
    \item Escribir usando lógica proposicional y realizar la tabla de verdad
    \item Asumiendo que la oración es verdadera y hay una torta, ¿qué se puede concluir?
    \item Asumiendo que la oración es verdadera y no hay una torta, ¿qué se puede concluir?
    \item Suponiendo que la oración es mentira (es falsa), ¿se puede concluir algo?
\end{itemize}

\vspace{0.2cm}
{\color{gray}\hrule}
\vspace{0.4cm}

\textbf{Resolución:}

Definimos las proposiciones simples: \\
$a = \text{Es mi cumpleaños}$ \\
$b = \text{Hay torta}$

\begin{enumerate}[label=\textbf{\arabic*)}]
    \item \textbf{Escribir usando lógica proposicional y realizar la tabla de verdad} \\
    La oración se traduce como: $(a \vee b) \Rightarrow b$

    \begin{center}
    \begin{tabular}{|c|c|c|c|}
    \hline
    $a$ & $b$ & $(a \vee b)$ & $(a \vee b) \Rightarrow b$ \\
    \hline
    V & V & V & \textbf{V} \\
    V & F & V & \textbf{F} \\
    F & V & V & \textbf{V} \\
    F & F & F & \textbf{V} \\
    \hline
    \end{tabular}
    \end{center}

    \item \textbf{Oración verdadera y hay torta} \\
    $(a \vee b) \Rightarrow b = \text{True}$, $b = V$. Entonces la oración es True, $b$ es True y $a$ no tiene un valor concreto, \underline{puede ser True o False}.

    \item \textbf{Oración verdadera y no hay torta} \\
    $(a \vee b) \Rightarrow b = \text{True}$, $b = F$. Entonces necesariamente el antecedente $(a \vee b) = F$, lo que implica que \underline{$a = \text{False}$}.

    \item \textbf{Oración es mentira (falsa)} \\
    $(a \vee b) \Rightarrow b = \text{False}$. Para que una implicación sea falsa, necesariamente el consecuente debe ser falso ($b = F$) y el antecedente verdadero. Entonces $(a \vee b) = \text{True}$, por lo tanto \underline{$a = \text{True}$}.
\end{enumerate}

\vspace{0.5cm}
\hrule
\vspace{0.5cm}
\hypertarget{ej3}{}
\noindent \textbf{Ejercicio 3.} Usando reglas de equivalencia (conmutatividad, asociatividad, De Morgan, etc.) determinar si los siguientes pares de fórmulas son equivalencias. Indicar en cada paso qué regla se utilizó.

\begin{enumerate}[label=\alph*)]
    \item $(p \vee q) \wedge (p \vee r) \quad \text{y} \quad \neg p \rightarrow (q \wedge r)$
    \item $\neg(\neg p) \rightarrow (\neg(\neg p \wedge \neg q)) \quad \text{y} \quad p \rightarrow (p \vee q)$
    \item $((True \wedge p) \wedge (\neg p \vee False)) \rightarrow \neg(\neg p \vee q) \quad \text{y} \quad p \wedge \neg q$
    \item $(p \vee (\neg p \wedge q)) \quad \text{y} \quad \neg p \rightarrow q$
    \item $p \rightarrow (q \wedge \neg(q \rightarrow r)) \quad \text{y} \quad (\neg p \vee q) \wedge (\neg p \vee (q \wedge \neg r))$
\end{enumerate}

\vspace{0.2cm}
{\color{gray}\hrule}
\vspace{0.4cm}

\textbf{Resolución:}

\begin{enumerate}[label=\textbf{\alph*)}]
    
    \item \textbf{$(p \vee q) \wedge (p \vee r) \quad \text{y} \quad \neg p \rightarrow (q \wedge r)$} \\
    Partimos de la primera expresión y aplicamos la propiedad \textbf{Distributiva} a la inversa (sacando factor común $p$):
    \[ (p \vee q) \wedge (p \vee r) \equiv p \vee (q \wedge r) \]
    Luego, por la definición de \textbf{Implicación} ($x \rightarrow y \equiv \neg x \vee y$), sabemos que $p \equiv \neg(\neg p)$, entonces:
    \[ \neg(\neg p) \vee (q \wedge r) \equiv \neg p \rightarrow (q \wedge r) \]
    \underline{\textbf{Son equivalentes.}}

    \vspace{0.3cm}
    \item \textbf{$\neg(\neg p) \rightarrow (\neg(\neg p \wedge \neg q)) \quad \text{y} \quad p \rightarrow (p \vee q)$} \\
    Trabajamos el lado izquierdo. Primero aplicamos \textbf{Doble Negación} en el antecedente: $\neg(\neg p) \equiv p$. \\
    Luego aplicamos las \textbf{Leyes de De Morgan} y \textbf{Doble Negación} en el consecuente: 
    \[ \neg(\neg p \wedge \neg q) \equiv \neg(\neg p) \vee \neg(\neg q) \equiv p \vee q \]
    Reescribiendo la expresión completa nos queda exactamente el lado derecho:
    \[ p \rightarrow (p \vee q) \]
    \underline{\textbf{Son equivalentes.}}

    \vspace{0.3cm}
    \item \textbf{$((True \wedge p) \wedge (\neg p \vee False)) \rightarrow \neg(\neg p \vee q) \quad \text{y} \quad p \wedge \neg q$} \\
    Simplificamos el antecedente usando la propiedad de \textbf{Elemento Neutro}:
    \[ True \wedge p \equiv p \]
    \[ \neg p \vee False \equiv \neg p \]
    Nos queda la conjunción $(p \wedge \neg p)$, que por la ley de \textbf{Contradicción} es siempre $False$. \\
    La implicación queda: $False \rightarrow \neg(\neg p \vee q)$. \\
    Como un antecedente falso hace que cualquier implicación sea automáticamente verdadera, todo el lado izquierdo se reduce a $True$. \\
    Como $True \not\equiv p \wedge \neg q$ (ya que la derecha depende de los valores de $p$ y $q$), concluimos que: \\
    \underline{\textbf{No son equivalentes.}}

    \vspace{0.3cm}
    \item \textbf{$(p \vee (\neg p \wedge q)) \quad \text{y} \quad \neg p \rightarrow q$} \\
    Aplicamos propiedad \textbf{Distributiva} en el lado izquierdo:
    \[ p \vee (\neg p \wedge q) \equiv (p \vee \neg p) \wedge (p \vee q) \]
    Por ley del \textbf{Tercero Excluido}, $(p \vee \neg p) \equiv True$. \\
    Por \textbf{Elemento Neutro}, $True \wedge (p \vee q) \equiv p \vee q$. \\
    Ahora miramos el lado derecho y aplicamos la definición de \textbf{Implicación}:
    \[ \neg p \rightarrow q \equiv \neg(\neg p) \vee q \equiv p \vee q \]
    Ambas expresiones se reducen a lo mismo. \\
    \underline{\textbf{Son equivalentes.}}

    \vspace{0.3cm}
    \item \textbf{$p \rightarrow (q \wedge \neg(q \rightarrow r)) \quad \text{y} \quad (\neg p \vee q) \wedge (\neg p \vee (q \wedge \neg r))$} \\
    Trabajamos primero el lado izquierdo. Aplicamos definición de \textbf{Implicación} en el paréntesis más interno:
    \[ p \rightarrow (q \wedge \neg(\neg q \vee r)) \]
    Aplicamos \textbf{De Morgan} y \textbf{Doble Negación}:
    \[ p \rightarrow (q \wedge (q \wedge \neg r)) \]
    Por propiedad \textbf{Asociativa} e \textbf{Idempotencia} ($q \wedge q \equiv q$):
    \[ p \rightarrow (q \wedge \neg r) \]
    Aplicamos definición de \textbf{Implicación} en la principal:
    \[ \neg p \vee (q \wedge \neg r) \]
    Ahora comparamos con el lado derecho. Si al lado derecho le aplicamos la propiedad \textbf{Distributiva} a la inversa (sacando $\neg p$ como factor común de la disyunción), obtenemos exactamente lo mismo:
    \[ (\neg p \vee q) \wedge (\neg p \vee (q \wedge \neg r)) \equiv \neg p \vee (q \wedge (q \wedge \neg r)) \equiv \neg p \vee (q \wedge \neg r) \]
    \underline{\textbf{Son equivalentes.}}

\end{enumerate}

\vspace{0.5cm}
\hrule
\vspace{0.5cm}
\hypertarget{ej4}{}
\noindent \textbf{Ejercicio 4.} Determinar si las siguientes fórmulas son tautologías, contradicciones o contingencias.

\begin{multicols}{2}
\begin{enumerate}[label=\alph*)]
    \item $(p \vee \neg p)$
    \item $(p \wedge \neg p)$
    \item $((\neg p \vee q) \leftrightarrow (p \rightarrow q))$
    \item $((p \wedge q) \rightarrow p)$
    \item $((p \wedge (q \vee r)) \leftrightarrow ((p \wedge q) \vee (p \wedge r)))$
    \item $((p \rightarrow (q \rightarrow r)) \rightarrow ((p \rightarrow q) \rightarrow (p \rightarrow r)))$
\end{enumerate}
\end{multicols}

\vspace{0.2cm}
{\color{gray}\hrule}
\vspace{0.4cm}

\textbf{Resolución:}

\begin{enumerate}[label=\textbf{\alph*)}]

    \item \textbf{$(p \vee \neg p)$} \\
    Armamos la tabla de verdad clásica:
    \begin{center}
    \begin{tabular}{|c|c|c|}
    \hline
    $p$ & $\neg p$ & $(p \vee \neg p)$ \\
    \hline
    V & F & \textbf{V} \\
    F & V & \textbf{V} \\
    \hline
    \end{tabular}
    \end{center}
    Como todos los resultados son verdaderos independientemente del valor de $p$, es una \underline{\textbf{Tautología}}.

    \vspace{0.4cm}
    \item \textbf{$(p \wedge \neg p)$}
    \begin{center}
    \begin{tabular}{|c|c|c|}
    \hline
    $p$ & $\neg p$ & $(p \wedge \neg p)$ \\
    \hline
    V & F & \textbf{F} \\
    F & V & \textbf{F} \\
    \hline
    \end{tabular}
    \end{center}
    Como todos los resultados son falsos, es una \underline{\textbf{Contradicción}}.

    \vspace{0.4cm}
    \item \textbf{$((\neg p \vee q) \leftrightarrow (p \rightarrow q))$}
    \begin{center}
    \begin{tabular}{|c|c|c|c|c|}
    \hline
    $p$ & $q$ & $(\neg p \vee q)$ & $(p \rightarrow q)$ & Equivalencia ($\leftrightarrow$) \\
    \hline
    V & V & V & V & \textbf{V} \\
    V & F & F & F & \textbf{V} \\
    F & V & V & V & \textbf{V} \\
    F & F & V & V & \textbf{V} \\
    \hline
    \end{tabular}
    \end{center}
    Ambos lados tienen los mismos valores de verdad, por lo tanto la equivalencia es una \underline{\textbf{Tautología}}.

    \vspace{0.4cm}
    \item \textbf{$((p \wedge q) \rightarrow p)$}
    \begin{center}
    \begin{tabular}{|c|c|c|c|}
    \hline
    $p$ & $q$ & $(p \wedge q)$ & $((p \wedge q) \rightarrow p)$ \\
    \hline
    V & V & V & \textbf{V} \\
    V & F & F & \textbf{V} \\
    F & V & F & \textbf{V} \\
    F & F & F & \textbf{V} \\
    \hline
    \end{tabular}
    \end{center}
    Como no hay ningún caso donde el antecedente sea verdadero y el consecuente falso, la implicación siempre se cumple. Es una \underline{\textbf{Tautología}}.

    \vspace{0.4cm}
    \item \textbf{$((p \wedge (q \vee r)) \leftrightarrow ((p \wedge q) \vee (p \wedge r)))$} \\
    Esta es la propiedad distributiva de la conjunción sobre la disyunción. Evaluando las 8 combinaciones posibles para $p, q, r$, vemos que el bicondicional siempre se cumple:
    \begin{center}
    \begin{tabular}{|c|c|c|c|c|c|}
    \hline
    $p$ & $q$ & $r$ & Lado Izq. & Lado Der. & ($\leftrightarrow$) \\
    \hline
    V & V & V & V & V & \textbf{V} \\
    V & V & F & V & V & \textbf{V} \\
    V & F & V & V & V & \textbf{V} \\
    V & F & F & F & F & \textbf{V} \\
    F & V & V & F & F & \textbf{V} \\
    F & V & F & F & F & \textbf{V} \\
    F & F & V & F & F & \textbf{V} \\
    F & F & F & F & F & \textbf{V} \\
    \hline
    \end{tabular}
    \end{center}
    Es una \underline{\textbf{Tautología}}.

    \vspace{0.4cm}
    \item \textbf{$((p \rightarrow (q \rightarrow r)) \rightarrow ((p \rightarrow q) \rightarrow (p \rightarrow r)))$} \\
    Nuevamente, al evaluar las 8 combinaciones posibles, el resultado de la implicación principal nunca tiene un antecedente $V$ con consecuente $F$.
    \begin{center}
    \begin{tabular}{|c|c|c|c|c|c|}
    \hline
    $p$ & $q$ & $r$ & Lado Izq. & Lado Der. & ($\rightarrow$) \\
    \hline
    V & V & V & V & V & \textbf{V} \\
    V & V & F & F & F & \textbf{V} \\
    V & F & V & V & V & \textbf{V} \\
    V & F & F & V & V & \textbf{V} \\
    F & V & V & V & V & \textbf{V} \\
    F & V & F & V & V & \textbf{V} \\
    F & F & V & V & V & \textbf{V} \\
    F & F & F & V & V & \textbf{V} \\
    \hline
    \end{tabular}
    \end{center}
    Es una \underline{\textbf{Tautología}}.

\end{enumerate}

\vspace{0.5cm}
\hrule
\vspace{0.5cm}
\hypertarget{ej5}{}
\noindent \textbf{Ejercicio 5.} Dadas las proposiciones lógicas $\alpha$ y $\beta$, se dice que $\alpha$ es más fuerte que $\beta$ si y sólo si $\alpha \rightarrow \beta$ es una tautología. En este caso, también decimos que $\beta$ es más débil que $\alpha$. Determinar la relación de fuerza de los siguientes pares de fórmulas:

\begin{multicols}{2}
\begin{enumerate}[label=\alph*)]
    \item True, False
    \item $(p \wedge q)$, $(p \vee q)$
    \item $p$, $(p \wedge q)$
    \item $p$, $(p \vee q)$
    \item $p$, $q$
    \item $p$, $(p \rightarrow q)$
\end{enumerate}
\end{multicols}

\noindent ¿Cuál es la proposición más fuerte y cuál la más débil de las que aparecen en este ejercicio?

\vspace{0.2cm}
{\color{gray}\hrule}
\vspace{0.4cm}

\textbf{Resolución:}

Para determinar la relación de fuerza, debemos evaluar la implicación en ambos sentidos: $\alpha \rightarrow \beta$ y $\beta \rightarrow \alpha$. Si una de estas implicaciones resulta en una Tautología (todos sus valores verdaderos), entonces su antecedente es la proposición más fuerte.

\begin{enumerate}[label=\textbf{\alph*)}]

    \item \textbf{True, False} \\
    Llamemos $\alpha = True$ y $\beta = False$. Evaluamos las implicaciones directamente por definición:
    \begin{itemize}
        \item $\alpha \rightarrow \beta \equiv True \rightarrow False \equiv \text{False}$ (No es tautología)
        \item $\beta \rightarrow \alpha \equiv False \rightarrow True \equiv \text{True}$ (Es tautología)
    \end{itemize}
    Como el antecedente que genera la tautología es el False, concluimos que \underline{\textbf{$False$ es más fuerte que $True$}}.

    \vspace{0.4cm}
    \item \textbf{$(p \wedge q)$, $(p \vee q)$} \\
    Llamemos $\alpha = (p \wedge q)$ y $\beta = (p \vee q)$. Armamos la tabla de verdad evaluando el ida ($\rightarrow$) y la vuelta ($\leftarrow$):
    \begin{center}
    \begin{tabular}{|c|c|c|c|c|c|}
    \hline
    $p$ & $q$ & $\alpha : (p \wedge q)$ & $\beta : (p \vee q)$ & $\alpha \rightarrow \beta$ & $\beta \rightarrow \alpha$ \\
    \hline
    V & V & V & V & \textbf{V} & V \\
    V & F & F & V & \textbf{V} & F \\
    F & V & F & V & \textbf{V} & F \\
    F & F & F & F & \textbf{V} & V \\
    \hline
    \end{tabular}
    \end{center}
    Vemos que $\alpha \rightarrow \beta$ es una tautología. Por lo tanto, \underline{\textbf{$(p \wedge q)$ es más fuerte que $(p \vee q)$}}.

    \vspace{0.4cm}
    \item \textbf{$p$, $(p \wedge q)$} \\
    Llamemos $\alpha = p$ y $\beta = (p \wedge q)$:
    \begin{center}
    \begin{tabular}{|c|c|c|c|c|}
    \hline
    $p$ & $q$ & $\beta : (p \wedge q)$ & $\alpha \rightarrow \beta$ & $\beta \rightarrow \alpha$ \\
    \hline
    V & V & V & V & \textbf{V} \\
    V & F & F & F & \textbf{V} \\
    F & V & F & V & \textbf{V} \\
    F & F & F & V & \textbf{V} \\
    \hline
    \end{tabular}
    \end{center}
    La tautología se da en $\beta \rightarrow \alpha$. Por lo tanto, \underline{\textbf{$(p \wedge q)$ es más fuerte que $p$}}.

    \vspace{0.4cm}
    \item \textbf{$p$, $(p \vee q)$} \\
    Llamemos $\alpha = p$ y $\beta = (p \vee q)$:
    \begin{center}
    \begin{tabular}{|c|c|c|c|c|}
    \hline
    $p$ & $q$ & $\beta : (p \vee q)$ & $\alpha \rightarrow \beta$ & $\beta \rightarrow \alpha$ \\
    \hline
    V & V & V & \textbf{V} & V \\
    V & F & V & \textbf{V} & V \\
    F & V & V & \textbf{V} & F \\
    F & F & F & \textbf{V} & V \\
    \hline
    \end{tabular}
    \end{center}
    La tautología se da en $\alpha \rightarrow \beta$. Por lo tanto, \underline{\textbf{$p$ es más fuerte que $(p \vee q)$}}.

    \vspace{0.4cm}
    \item \textbf{$p$, $q$} \\
    Llamemos $\alpha = p$ y $\beta = q$:
    \begin{center}
    \begin{tabular}{|c|c|c|c|}
    \hline
    $p$ & $q$ & $p \rightarrow q$ & $q \rightarrow p$ \\
    \hline
    V & V & V & V \\
    V & F & F & V \\
    F & V & V & F \\
    F & F & V & V \\
    \hline
    \end{tabular}
    \end{center}
    Ninguna de las dos columnas genera una tautología. Por lo tanto, \underline{\textbf{No hay relación de fuerza}} (son proposiciones incomparables).

    \vspace{0.4cm}
    \item \textbf{$p$, $(p \rightarrow q)$} \\
    Llamemos $\alpha = p$ y $\beta = (p \rightarrow q)$:
    \begin{center}
    \begin{tabular}{|c|c|c|c|c|}
    \hline
    $p$ & $q$ & $\beta : (p \rightarrow q)$ & $\alpha \rightarrow \beta$ & $\beta \rightarrow \alpha$ \\
    \hline
    V & V & V & V & V \\
    V & F & F & F & V \\
    F & V & V & V & F \\
    F & F & V & V & F \\
    \hline
    \end{tabular}
    \end{center}
    Nuevamente, ninguna de las dos implicaciones resulta en una tautología completa. Por lo tanto, \underline{\textbf{No hay relación de fuerza}}.

\end{enumerate}

\vspace{0.4cm}
\textbf{Pregunta final: ¿Cuál es la proposición más fuerte y cuál la más débil?}

\begin{itemize}
    \item La \textbf{más fuerte de todas} es \textbf{$False$}. Una contradicción implica lógicamente a cualquier otra proposición (el principio de \textit{ex falso quodlibet}). Al ser el antecedente falso, la implicación siempre es verdadera.
    \item La \textbf{más débil de todas} es \textbf{$True$}. Una tautología es implicada por cualquier proposición. Al ser el consecuente verdadero, no importa el valor del antecedente, la implicación siempre se cumple.
\end{itemize}

\vspace{0.5cm}
\hrule
\vspace{0.5cm}

\section*{1.2. Lógica trivaluada}

\hypertarget{ej6}{}
\noindent \textbf{Ejercicio 6.} Asumiendo que el valor de verdad de $b$ y $c$ es verdadero, el de $a$ es falso y el de $x$ e $y$ es indefinido, indicar cuáles de los operadores deben ser operadores "luego" para que la expresión no se indefina nunca:
\begin{multicols}{2}
\begin{enumerate}[label=\alph*)]
    \item $(\neg x \vee b)$
    \item $((c \vee (y \wedge a)) \vee b)$
    \item $\neg(c \vee y)$
    \item $(\neg(c \vee y) \leftrightarrow (\neg c \wedge \neg y))$
    \item $((c \vee y) \wedge (a \vee b))$
    \item $(((c \vee y) \wedge (a \vee b)) \leftrightarrow (c \vee (y \wedge a) \vee b))$
    \item $(\neg c \wedge \neg y)$
\end{enumerate}
\end{multicols}

\vspace{0.2cm}
{\color{gray}\hrule}
\vspace{0.4cm}

\textbf{Resolución:}

\textbf{Repaso teórico:} En lógica trivaluada (y en la programación en general), los operadores "luego" emplean \textbf{evaluación de cortocircuito}. Al evaluar estrictamente de izquierda a derecha, se detienen y devuelven un valor apenas el resultado ya está garantizado, evitando evaluar términos conflictivos o indefinidos ($\bot$):
\begin{itemize}
    \item \textbf{Disyunción luego ($\vee_L$):} Si el lado izquierdo es $V$, devuelve $V$ sin mirar el lado derecho ($V \vee_L \bot \equiv V$).
    \item \textbf{Conjunción luego ($\wedge_L$):} Si el lado izquierdo es $F$, devuelve $F$ sin mirar el lado derecho ($F \wedge_L \bot \equiv F$).
\end{itemize}

Valores dados: $b=V$, $c=V$, $a=F$, $x=\bot$, $y=\bot$. Analizamos caso por caso:

\begin{enumerate}[label=\textbf{\alph*)}]

    \vspace{0.3cm}
    \item \textbf{$(\neg x \vee b)$} \\
    Reemplazamos: $(\neg \bot \vee V)$. Al evaluar de izquierda a derecha, lo primero que el programa intenta hacer es resolver $\neg x$. Como $x$ es indefinido, la expresión falla de inmediato en ese punto. \\
    \underline{\textbf{Respuesta:}} No se puede evitar que se indefina (falla siempre). Ni siquiera un operador "luego" lo salva, porque el término indefinido está a la izquierda.

    \vspace{0.4cm}
    \item \textbf{$((c \vee (y \wedge a)) \vee b)$} \\
    Reemplazamos: $((V \vee (\bot \wedge F)) \vee V)$. Evaluamos de izquierda a derecha. En el primer paréntesis leemos la $c$ que vale $V$. Si el $\vee$ que le sigue es "luego", hace cortocircuito y todo el bloque $((V \vee_L (\bot \wedge F))$ se transforma en $V$ sin tocar la $y$. \\
    \underline{\textbf{Respuesta:}} El $\vee$ entre la $c$ y $(y \wedge a)$ debe ser un operador "luego".

    \vspace{0.4cm}
    \item \textbf{$\neg(c \vee y)$} \\
    Reemplazamos: $\neg(V \vee \bot)$. Para evitar que se evalúe la $y$ indefinida, necesitamos que la disyunción dentro del paréntesis haga cortocircuito al leer la $V$ inicial. \\
    \underline{\textbf{Respuesta:}} El operador $\vee$ debe ser un operador "luego".

    \vspace{0.4cm}
    \item \textbf{$(\neg(c \vee y) \leftrightarrow (\neg c \wedge \neg y))$} \\
    Analizamos ambos lados del bicondicional:
    \begin{itemize}
        \item \textbf{Lado izquierdo:} $\neg(V \vee \bot)$. Igual que en el inciso c), necesitamos que sea $\vee_L$ para que evalúe a $\neg(V) \equiv F$.
        \item \textbf{Lado derecho:} $(\neg V \wedge \neg \bot) \equiv (F \wedge \neg \bot)$. Al leer el primer valor, obtenemos un $F$. Si la conjunción es "luego", hace cortocircuito y devuelve $F$ de inmediato, evitando el $\neg \bot$.
    \end{itemize}
    \underline{\textbf{Respuesta:}} El $\vee$ del lado izquierdo y el $\wedge$ del lado derecho deben ser operadores "luego".

    \vspace{0.4cm}
    \item \textbf{$((c \vee y) \wedge (a \vee b))$} \\
    Reemplazamos: $((V \vee \bot) \wedge (F \vee V))$. El único peligro de indefinición está en el primer paréntesis. Como arranca con $V$, un $\vee_L$ lo hace valer $V$ automáticamente por cortocircuito. El segundo paréntesis no tiene indefiniciones, así que no importa qué operadores tenga. \\
    \underline{\textbf{Respuesta:}} El $\vee$ entre $c$ e $y$ debe ser un operador "luego".

    \vspace{0.4cm}
    \item \textbf{$(((c \vee y) \wedge (a \vee b)) \leftrightarrow (c \vee (y \wedge a) \vee b))$} \\
    Analizamos por partes:
    \begin{itemize}
        \item \textbf{Lado izquierdo:} Igual que en el inciso e), el primer $\vee$ (entre $c$ e $y$) debe ser $\vee_L$. Todo el bloque izquierdo terminará valiendo $V$.
        \item \textbf{Lado derecho:} $(V \vee (\bot \wedge F) \vee V)$. Evaluando de izquierda a derecha, la proposición arranca con $V$. Si el primer $\vee$ (el que le sigue a la $c$) es un $\vee_L$, hace cortocircuito y omite evaluar todo el resto de la expresión, devolviendo $V$.
    \end{itemize}
    \underline{\textbf{Respuesta:}} El primer $\vee$ del lado izquierdo y el primer $\vee$ del lado derecho deben ser operadores "luego".

    \vspace{0.4cm}
    \item \textbf{$(\neg c \wedge \neg y)$} \\
    Reemplazamos: $(\neg V \wedge \neg \bot) \equiv (F \wedge \neg \bot)$. Evaluando de izquierda a derecha, ya sabemos que el primer valor es falso ($F$). Si usamos una conjunción de cortocircuito ($\wedge_L$), la expresión entera devuelve $F$ sin intentar evaluar el $\neg y$. \\
    \underline{\textbf{Respuesta:}} El operador $\wedge$ debe ser un operador "luego".

\end{enumerate}

\vspace{0.5cm}
\hrule
\vspace{0.5cm}
\hypertarget{ej7}{}
\noindent \textbf{Ejercicio 7.} Sean $p, q$ y $r$ tres variables de las que se sabe que:
\begin{itemize}
    \item $p$ y $q$ nunca están indefinidas,
    \item $r$ se indefine sii $q$ es verdadera
\end{itemize}
Proponer una fórmula que nunca se indefina, utilizando siempre las tres variables y que sea verdadera si y solo si se cumple que:

\begin{multicols}{2}
\begin{enumerate}[label=\alph*)]
    \item Al menos una es verdadera.
    \item Ninguna es verdadera.
    \item Exactamente una de las tres es verdadera.
    \item Sólo $p$ y $q$ son verdaderas.
    \item No todas al mismo tiempo son verdaderas.
    \item $r$ es verdadera.
\end{enumerate}
\end{multicols}

\vspace{0.2cm}
{\color{gray}\hrule}
\vspace{0.4cm}

\textbf{Resolución:}

El dato clave del ejercicio es: **Si $q=V$, entonces $r=\bot$**. 
Esto significa que nunca podemos evaluar $r$ si $q$ es verdadero. Debemos proteger siempre a $r$ usando operadores de cortocircuito, típicamente poniéndola a la derecha de un $(\neg q \wedge_L \dots)$ o un $(q \vee_L \dots)$. Además, el enunciado exige que las fórmulas utilicen explícitamente las tres variables.

\begin{enumerate}[label=\textbf{\alph*)}]

    \item \textbf{Al menos una es verdadera.} \\
    Fórmula: $p \vee q \vee_L r$ \\
    \textit{Justificación:} Evaluando de izquierda a derecha, si $q$ es $V$, el operador $\vee_L$ hace cortocircuito y devuelve $V$ sin mirar $r$ (lo cual es correcto porque ya hay al menos una verdadera). Si $q$ es $F$, es seguro evaluar $r$, y la proposición completa resultará verdadera si $p$ o $r$ lo son.

    \vspace{0.4cm}
    \item \textbf{Ninguna es verdadera.} \\
    Fórmula: $\neg p \wedge \neg q \wedge_L \neg r$ \\
    \textit{Justificación:} Si $q$ es $V$, entonces $\neg q$ es $F$. El operador $\wedge_L$ hace cortocircuito y devuelve $F$ automáticamente, protegiendo a $r$. Esto es correcto porque si $q$ es $V$, es imposible que "ninguna sea verdadera". Si $q$ es $F$, entonces $\neg q$ es $V$, y se procede a evaluar $r$ de forma segura.

    \vspace{0.4cm}
    \item \textbf{Exactamente una de las tres es verdadera.} \\
    Fórmula: $(q \wedge \neg p) \vee (\neg q \wedge_L ((p \wedge \neg r) \vee (\neg p \wedge r)))$ \\
    \textit{Justificación:} Dividimos en dos casos lógicos:
    \begin{itemize}
        \item Si $q=V$: Sabemos que $r=\bot$ (no es verdadera). Para que "exactamente una" sea verdadera, $p$ debe ser $F$. El primer paréntesis evalúa a $\neg p$. El segundo paréntesis hace cortocircuito en $\neg q = F$ y no evalúa $r$.
        \item Si $q=F$: El primer paréntesis es $F$. El segundo pasa a evaluar la disyunción exclusiva entre $p$ y $r$, verificando que exactamente una de las dos restantes sea verdadera.
    \end{itemize}

    \vspace{0.4cm}
    \item \textbf{Sólo $p$ y $q$ son verdaderas.} \\
    Fórmula: $p \wedge q \wedge (q \vee_L r)$ \\
    \textit{Justificación:} Si $q=V$, el paréntesis $(q \vee_L r)$ hace cortocircuito y da $V$. La fórmula queda evaluando $p \wedge V \wedge V \equiv p$. Si $q=F$, el paréntesis evalúa a $r$, pero como tenemos $q$ en la conjunción principal, toda la expresión da $F$ automáticamente, lo cual es correcto ya que exigíamos que $q$ sea verdadera.

    \vspace{0.4cm}
    \item \textbf{No todas al mismo tiempo son verdaderas.} \\
    Fórmula: $\neg (p \wedge q \wedge (\neg q \wedge_L r))$ \\
    \textit{Justificación:} Analicemos la frase: como si $q=V \Rightarrow r=\bot$, es imposible que $q$ y $r$ sean verdaderas al mismo tiempo en cualquier escenario. Por lo tanto, "no todas son verdaderas a la vez" es una Tautología (siempre es V). Construimos una fórmula que siempre da Falso para luego negarla: El término $(\neg q \wedge_L r)$ protege a $r$ de forma segura. Como está en conjunción con $q$, siempre da Falso. La negación exterior lo vuelve siempre Verdadero.

    \vspace{0.4cm}
    \item \textbf{$r$ es verdadera.} \\
    Fórmula: $(\neg q \wedge_L r) \wedge (p \vee \neg p)$ \\
    \textit{Justificación:} Si $q=V$, $r=\bot$, por lo que $r$ no es verdadera (es indefinida). La expresión debe dar $F$. El término $(\neg q \wedge_L r)$ hace cortocircuito en $F$ y lo logra. Si $q=F$, evalúa $r$ y devuelve su valor. Multiplicamos la expresión por la tautología $(p \vee \neg p)$ únicamente para cumplir con el requisito del enunciado de "utilizar siempre las tres variables" sin alterar el valor de verdad.

\end{enumerate}

\vspace{0.5cm}
\hrule
\vspace{0.5cm}

\section*{1.3. Lógica de primer orden}

\hypertarget{ej8}{}
\noindent \textbf{Ejercicio 8.} Determinar, para cada aparición de variables, si dicha aparición se encuentra libre o ligada. En caso de estar ligada, aclarar a qué cuantificador lo está. En los casos en que sea posible, proponer valores para las variables libres de modo tal que las expresiones sean verdaderas.

\begin{enumerate}[label=\alph*)]
    \item $(\forall x: \mathbb{Z}) (0 \leq x < n \rightarrow x + y = z)$
    \item $(\exists x: \mathbb{Z}) ((\forall y: \mathbb{Z}) (0 \leq x < n \wedge 0 < y < m \rightarrow x + y = z))$
    \item $(\forall j: \mathbb{Z}) (0 \leq j < 10 \rightarrow j < 0)$
    \item $(\forall j: \mathbb{Z}) (j \leq 0 \rightarrow P(j)) \wedge P(j)$
\end{enumerate}

\vspace{0.2cm}
{\color{gray}\hrule}
\vspace{0.4cm}

\textbf{Resolución:}

Recordemos la teoría básica: una variable está \textbf{ligada} si se encuentra bajo el alcance de un cuantificador (como $\forall$ o $\exists$) que la referencia explícitamente. Está \textbf{libre} si su valor depende del contexto exterior y no está atada a ningún cuantificador en la expresión.

\begin{enumerate}[label=\textbf{\alph*)}]

    \item \textbf{$(\forall x: \mathbb{Z}) (0 \leq x < n \rightarrow x + y = z)$} \\
    \textbf{Análisis de variables:}
    \begin{itemize}
        \item $x$: \textbf{Ligada} al cuantificador universal $(\forall x: \mathbb{Z})$.
        \item $n$, $y$, $z$: \textbf{Libres}.
    \end{itemize}
    \textbf{Valores para que sea verdadera:} Proponemos $n = 0$, con $y, z \in \mathbb{Z}$ arbitrarios. Al ser $n=0$, la condición del antecedente $0 \leq x < 0$ siempre es falsa. Como un antecedente falso hace que la implicación completa sea Verdadera (vacuidad), la expresión es Verdadera.

    \vspace{0.4cm}
    \item \textbf{$(\exists x: \mathbb{Z}) ((\forall y: \mathbb{Z}) (0 \leq x < n \wedge 0 < y < m \rightarrow x + y = z))$} \\
    \textbf{Análisis de variables:}
    \begin{itemize}
        \item $x$: \textbf{Ligada} al cuantificador existencial $(\exists x: \mathbb{Z})$.
        \item $y$: \textbf{Ligada} al cuantificador universal $(\forall y: \mathbb{Z})$.
        \item $n$, $m$, $z$: \textbf{Libres}.
    \end{itemize}
    \textbf{Valores para que sea verdadera:} Usamos la misma lógica del inciso anterior. Proponemos $m = 0$ (con $n, z \in \mathbb{Z}$ arbitrarios). De esta forma, el antecedente $0 < y < 0$ es absurdo (Falso), la implicación se vuelve trivialmente verdadera, y por lo tanto existe cualquier $x$ que la cumple.

    \vspace{0.4cm}
    \item \textbf{$(\forall j: \mathbb{Z}) (0 \leq j < 10 \rightarrow j < 0)$} \\
    \textbf{Análisis de variables:}
    \begin{itemize}
        \item $j$: \textbf{Ligada} al cuantificador universal $(\forall j: \mathbb{Z})$.
    \end{itemize}
    \textbf{Valores para que sea verdadera:} Esta expresión \textbf{no tiene variables libres} a las que asignarles valores (es una fórmula cerrada). Además, evalúa siempre a \textbf{Falsa}. (Por ejemplo, si instanciamos $j=1$, el antecedente $0 \leq 1 < 10$ es $V$, pero el consecuente $1 < 0$ es $F$, resultando en $V \rightarrow F \equiv F$). Por lo tanto, no es posible proponer valores.

    \vspace{0.4cm}
    \item \textbf{$(\forall j: \mathbb{Z}) (j \leq 0 \rightarrow P(j)) \wedge P(j)$} \\
    \textbf{Análisis de variables:} \\
    ¡Atención acá con el alcance de los paréntesis!
    \begin{itemize}
        \item $j$ (primeras dos apariciones): \textbf{Ligada} al cuantificador $(\forall j: \mathbb{Z})$. El alcance termina al cerrar el paréntesis de la implicación.
        \item $j$ (última aparición en el último $P(j)$): \textbf{Libre}. Queda afuera del cuantificador.
        \item $P$: \textbf{Libre} (es un símbolo de predicado libre que debemos instanciar).
    \end{itemize}
    \textbf{Valores para que sea verdadera:} Proponemos que el predicado libre sea una tautología, por ejemplo $P(x) = True$ para todo $x$, y que la variable libre sea $j = 1$. Así, la implicación $(\forall j) (j \leq 0 \rightarrow True)$ es verdadera, y el término suelto $P(1)$ también es verdadero. Resultando en $True \wedge True \equiv True$.

\end{enumerate}

\vspace{0.5cm}
\hrule
\vspace{0.5cm}
\hypertarget{ej9}{}
\noindent \textbf{Ejercicio 9.} Sea $P(x:\mathbb{Z})$ y $Q(x:\mathbb{Z})$ dos predicados cualquiera. Explicar cuál es el error de traducción a fórmulas de los siguientes enunciados. Dar un ejemplo en el cuál sucede el problema y luego corregirlo.

\begin{enumerate}[label=\alph*)]
    \item "Todos los naturales menores a 10 cumple P" \\
    $(\forall i:\mathbb{Z})(0 \leq i < 10 \wedge P(i))$
    \item "Algún natural menor a 10 cumple P" \\
    $(\exists i:\mathbb{Z})(0 \leq i < 10 \rightarrow P(i))$
    \item "Todos los naturales menores a 10 que cumplen P, cumplen Q" \\
    $(\forall x:\mathbb{Z})(0 \leq x < 10 \rightarrow (P(x) \wedge Q(x)))$
    \item "No hay ningún natural menor a 10 que cumpla P y Q" \\
    $\neg (\exists x:\mathbb{Z})(0 \leq x < 10 \wedge P(x)) \wedge \neg (\exists x:\mathbb{Z})(0 \leq x < 10 \wedge Q(x))$
\end{enumerate}

\vspace{0.2cm}
{\color{gray}\hrule}
\vspace{0.4cm}

\textbf{Resolución:}

\begin{enumerate}[label=\textbf{\alph*)}]

    \item \textbf{Error:} Se utilizó una conjunción ($\wedge$) para filtrar el dominio dentro de un cuantificador universal ($\forall$). \\
    \textbf{Ejemplo de fallo:} Si evaluamos la fórmula para $i=15$, la primera parte de la conjunción resulta falsa ($0 \leq 15 < 10 \equiv Falso$). Como es una conjunción, toda la expresión para $i=15$ es Falsa. Y como el cuantificador exige que \textit{para todo} entero la expresión sea Verdadera, toda la proposición colapsa y se vuelve Falsa, sin importar si los menores a 10 realmente cumplían $P$ o no. \\
    \textbf{Corrección:} Para filtrar en un $\forall$, debemos usar la implicación ($\rightarrow$). Así, los números fuera del rango hacen falso al antecedente y la implicación da Verdadero (por vacuidad), dejándonos evaluar $P$ solo donde nos importa.
    \[ (\forall i:\mathbb{Z})(0 \leq i < 10 \rightarrow P(i)) \]

    \vspace{0.4cm}
    \item \textbf{Error:} Se utilizó una implicación ($\rightarrow$) dentro de un cuantificador existencial ($\exists$). \\
    \textbf{Ejemplo de fallo:} Si evaluamos para $i=15$, el antecedente de la implicación es falso ($0 \leq 15 < 10 \equiv Falso$). Por definición de implicación, un antecedente falso hace que la expresión completa sea Verdadera. Por lo tanto, el $\exists$ encuentra a $i=15$ (o a cualquier otro número fuera del rango) como un caso de éxito y la fórmula devuelve Verdadero, afirmando que "existe un natural menor a 10 que cumple P" basándose en un número que ni siquiera está en el rango. \\
    \textbf{Corrección:} Para exigir condiciones simultáneas en un $\exists$, debemos usar la conjunción ($\wedge$).
    \[ (\exists i:\mathbb{Z})(0 \leq i < 10 \wedge P(i)) \]

    \vspace{0.4cm}
    \item \textbf{Error:} La fórmula propuesta significa "Todos los naturales menores a 10 cumplen P \textit{y además} cumplen Q". El enunciado original no exige que todos cumplan P; dice que \textit{si} lo cumplen, entonces también deben cumplir Q (es una condición). \\
    \textbf{Ejemplo de fallo:} Supongamos que el número 3 no cumple P. La fórmula propuesta evaluará a Falso para $x=3$ porque $(F \wedge Q(3)) \equiv F$, y por ende el $\forall$ dará Falso. Sin embargo, según el enunciado original en lenguaje natural, que el 3 no cumpla P no es un problema (simplemente no entra en el grupo de los que deben cumplir Q). \\
    \textbf{Corrección:} Agrupar la pertenencia al rango y el cumplimiento de P como antecedente de la implicación hacia Q.
    \[ (\forall x:\mathbb{Z})((0 \leq x < 10 \wedge P(x)) \rightarrow Q(x)) \]

    \vspace{0.4cm}
    \item \textbf{Error:} La fórmula distribuyó la negación existencial, lo que cambió drásticamente el significado. La fórmula propuesta afirma: "No hay ningún natural menor a 10 que cumpla P, \textit{y además} no hay ningún natural menor a 10 que cumpla Q" (es decir, ninguno cumple ninguna de las dos). El enunciado original prohíbe que alguien cumpla \textit{ambas al mismo tiempo}. \\
    \textbf{Ejemplo de fallo:} Supongamos que $x=2$ cumple P pero no Q, y $x=3$ cumple Q pero no P. El enunciado original sería Verdadero (porque nadie cumple ambas a la vez). Sin embargo, la fórmula propuesta será Falsa, porque la primera parte $\neg(\exists x \dots P(x))$ evalúa a Falso (¡ya que el 2 sí cumple P!). \\
    \textbf{Corrección:} Negar la existencia de un elemento que cumpla todas las condiciones en simultáneo mediante una gran conjunción.
    \[ \neg (\exists x:\mathbb{Z})(0 \leq x < 10 \wedge P(x) \wedge Q(x)) \]

\end{enumerate}

\vspace{0.5cm}
\hrule
\vspace{0.5cm}
\vspace{1cm}
\hypertarget{ej10}{}
\noindent \textbf{Ejercicio 10.} Sean $P(x:\mathbb{Z})$ y $Q(x:\mathbb{Z})$ dos predicados cualesquiera que nunca se indefinen. Escribir el predicado asociado a cada uno de los siguientes enunciados:

\begin{enumerate}[label=\alph*)]
    \item “Existe un único número natural menor a 10 que cumple P”
    \item “Existen al menos dos números naturales menores a 10 que cumplen P”
    \item “Existen exactamente dos números naturales menores a 10 que cumplen P”
    \item “Todos los enteros pares que cumplen P, no cumplen Q”
    \item “Si un entero cumple P y es impar, no cumple Q”
    \item “Todos los enteros pares cumplen P, y todos los enteros impares que no cumplen P cumplen Q”
    \item “Si hay un número natural menor a 10 que no cumple P entonces ninguno natural menor a 10 cumple Q; y si todos los naturales menores a 10 cumplen P entonces hay al menos dos naturales menores a 10 que cumplen Q”
\end{enumerate}

\vspace{0.2cm}
{\color{gray}\hrule}
\vspace{0.4cm}

\textbf{Resolución:}

\textit{Nota: Siguiendo la convención de la materia vista en el Ejercicio 9, traducimos "número natural menor a 10" usando el dominio $\mathbb{Z}$ con el rango $0 \leq x < 10$.}

\begin{enumerate}[label=\textbf{\alph*)}]

    \item \textbf{“Existe un único número natural menor a 10 que cumple P“} \\
    Para asegurar que es "único", primero afirmamos que existe al menos uno que cumple. Luego, garantizamos que cualquier otro número que también cumpla, debe ser obligatoriamente el mismo que ya encontramos.
    \[ (\exists x: \mathbb{Z})(0 \leq x < 10 \wedge P(x) \wedge (\forall y: \mathbb{Z})(0 \leq y < 10 \wedge P(y) \rightarrow x = y)) \]

    \vspace{0.4cm}
    \item \textbf{“Existen al menos dos números naturales menores a 10 que cumplen P”} \\
    Pedimos la existencia de dos variables distintas. El "al menos" significa que no nos importa si hay 3, 4 o más; con encontrar dos diferentes es suficiente. La clave es exigir que $x \neq y$.
    \[ (\exists x, y: \mathbb{Z})(0 \leq x < 10 \wedge 0 \leq y < 10 \wedge x \neq y \wedge P(x) \wedge P(y)) \]

    \vspace{0.4cm}
    \item \textbf{“Existen exactamente dos números naturales menores a 10 que cumplen P”} \\
    Es la combinación de las dos lógicas anteriores. Afirmamos que existen dos distintos (al menos dos), y luego usamos el $\forall$ para "cerrar la puerta" y asegurar que cualquier otro elemento $z$ que intente cumplir $P$, tiene que ser idéntico a $x$ o a $y$.
    \begin{multline*}
    (\exists x, y: \mathbb{Z})(0 \leq x < 10 \wedge 0 \leq y < 10 \wedge x \neq y \wedge P(x) \wedge P(y) \wedge \\
    (\forall z: \mathbb{Z})(0 \leq z < 10 \wedge P(z) \rightarrow (z = x \vee z = y)))
    \end{multline*}

    \vspace{0.4cm}
    \item \textbf{“Todos los enteros pares que cumplen P, no cumplen Q”} \\
    La paridad de un entero se expresa analizando su resto en la división por 2 ($x \bmod 2 = 0$). Como habla de "todos", usamos un $\forall$ con una implicación.
    \[ (\forall x: \mathbb{Z})(x \bmod 2 = 0 \wedge P(x) \rightarrow \neg Q(x)) \]

    \vspace{0.4cm}
    \item \textbf{“Si un entero cumple P y es impar, no cumple Q”} \\
    Es un caso análogo al anterior. La frase "si un entero..." es una forma condicional que aplica universalmente. Que sea impar se escribe como $x \bmod 2 \neq 0$.
    \[ (\forall x: \mathbb{Z})(P(x) \wedge x \bmod 2 \neq 0 \rightarrow \neg Q(x)) \]

    \vspace{0.4cm}
    \item \textbf{"Todos los enteros pares cumplen P, y todos los enteros impares que no cumplen P cumplen Q"} \\
    Son dos afirmaciones universales distintas unidas por una gran conjunción. Podemos dividirlo claramente en dos cuantificadores $\forall$ separados.
    \[ (\forall x: \mathbb{Z})(x \bmod 2 = 0 \rightarrow P(x)) \wedge (\forall y: \mathbb{Z})(y \bmod 2 \neq 0 \wedge \neg P(y) \rightarrow Q(y)) \]

    \vspace{0.4cm}
    \item \textbf{“Si hay un número natural menor a 10 que no cumple P entonces ninguno natural menor a 10 cumple Q; y si todos los naturales menores a 10 cumplen P entonces hay al menos dos naturales menores a 10 que cumplen Q”} \\
    Este es un mega-enunciado formado por dos grandes implicaciones unidas por un "y" lógico ($\wedge$). Traducimos cada bloque por separado: \\
    \textit{Bloque 1:} $(\exists x: \mathbb{Z})(0 \leq x < 10 \wedge \neg P(x)) \rightarrow (\forall y: \mathbb{Z})(0 \leq y < 10 \rightarrow \neg Q(y))$ \\
    \textit{Bloque 2:} $(\forall z: \mathbb{Z})(0 \leq z < 10 \rightarrow P(z)) \rightarrow (\exists v, w: \mathbb{Z})(0 \leq v < 10 \wedge 0 \leq w < 10 \wedge v \neq w \wedge Q(v) \wedge Q(w))$ \\
    Juntando todo queda:
    \begin{multline*}
    \Big( (\exists x: \mathbb{Z})(0 \leq x < 10 \wedge \neg P(x)) \rightarrow (\forall y: \mathbb{Z})(0 \leq y < 10 \rightarrow \neg Q(y)) \Big) \wedge \\
    \Big( (\forall z: \mathbb{Z})(0 \leq z < 10 \rightarrow P(z)) \rightarrow \\
    (\exists v, w: \mathbb{Z})(0 \leq v < 10 \wedge 0 \leq w < 10 \wedge v \neq w \wedge Q(v) \wedge Q(w)) \Big)
    \end{multline*}

\end{enumerate}

\vspace{0.5cm}
\hrule
\vspace{0.5cm}
\hypertarget{ej11}{}
\noindent \textbf{Ejercicio 11.} Sean $P(x:\mathbb{Z})$ y $Q(x:\mathbb{Z}, y:\mathbb{Z})$ dos predicados cualesquiera que nunca se indefinen. Escribir el predicado asociado a cada uno de los siguientes enunciados:

\begin{enumerate}[label=\alph*)]
    \item “Si un entero cumple P, entonces existe un entero distinto tal que juntos cumplen Q“
    \item “Existe un par de enteros tal que cumplen Q y ninguno de ambos cumple P“
    \item “Si un par de enteros cumplen Q, entonces al menos uno de ellos cumple P“
    \item “Si un entero cumple P, no existe ningún entero tal que juntos cumplan Q“
\end{enumerate}

\vspace{0.2cm}
{\color{gray}\hrule}
\vspace{0.4cm}

\textbf{Resolución:}

\begin{enumerate}[label=\textbf{\alph*)}]

    \item \textbf{“Si un entero cumple P, entonces existe un entero distinto tal que juntos cumplen Q“} \\
    La frase "Si un entero..." establece una condición general, por lo que usamos un cuantificador universal ($\forall$). Para decir que es "distinto", exigimos explícitamente que la nueva variable sea diferente a la primera ($x \neq y$).
    \[ (\forall x: \mathbb{Z})(P(x) \rightarrow (\exists y: \mathbb{Z})(x \neq y \wedge Q(x, y))) \]

    \vspace{0.4cm}
    \item \textbf{“Existe un par de enteros tal que cumplen Q y ninguno de ambos cumple P“} \\
    Acá buscamos la existencia de dos variables. La palabra "par" sugiere fuertemente que son dos elementos distintos, por lo que es una buena práctica (y más riguroso) agregar $x \neq y$. Que "ninguno" cumpla $P$ significa que tanto $x$ como $y$ hacen falso al predicado $P$.
    \[ (\exists x, y: \mathbb{Z})(x \neq y \wedge Q(x, y) \wedge \neg P(x) \wedge \neg P(y)) \]

    \vspace{0.4cm}
    \item \textbf{“Si un par de enteros cumplen Q, entonces al menos uno de ellos cumple P“} \\
    Nuevamente, el "Si un par..." actúa como una regla general, por lo que usamos un $\forall$. Si la condición (cumplir $Q$) se da, entonces $P(x)$ o $P(y)$ debe ser verdadero (disyunción).
    \[ (\forall x, y: \mathbb{Z})(x \neq y \wedge Q(x, y) \rightarrow (P(x) \vee P(y))) \]
    \textit{Nota: Dependiendo de cuán estricto sea el docente con la palabra "par", el $x \neq y$ en el antecedente puede omitirse, pero ponerlo demuestra que entendemos que hablamos de dos enteros diferentes.}

    \vspace{0.4cm}
    \item \textbf{“Si un entero cumple P, no existe ningún entero tal que juntos cumplan Q“} \\
    Arrancamos con el $\forall$ general para el primer entero. Luego, en el consecuente, negamos literalmente la existencia de un segundo entero que cumpla la relación $Q$ con el primero.
    \[ (\forall x: \mathbb{Z})(P(x) \rightarrow \neg (\exists y: \mathbb{Z}) Q(x, y)) \]
    \textit{Alternativa equivalente (pasando la negación hacia adentro):} 
    \[ (\forall x: \mathbb{Z})(P(x) \rightarrow (\forall y: \mathbb{Z}) \neg Q(x, y)) \]

\end{enumerate}

\vspace{0.5cm}
\hrule
\vspace{0.5cm}
\hypertarget{ej12}{}
\noindent \textbf{Ejercicio 12.} Sean $P(x:\mathbb{Z})$ y $Q(x:\mathbb{Z})$ dos predicados cualesquiera que nunca se indefinen y sean a, b y k enteros. Decidir en cada caso la relación de fuerza entre las dos fórmulas:

\begin{enumerate}[label=\alph*)]
    \item $P(3) \quad \text{y} \quad (\forall n: \mathbb{Z}) (0 \leq n < 5 \rightarrow P(n))$
    \item $P(3) \quad \text{y} \quad (\exists n: \mathbb{Z}) (0 \leq n < 5 \wedge P(n))$
    \item $(\forall n:\mathbb{Z})((0 \leq n < 10 \wedge P(n)) \rightarrow Q(n)) \quad \text{y} \quad (\forall n: \mathbb{Z}) (0 \leq n < 10 \rightarrow Q(n))$
    \item $(\exists n:\mathbb{Z})(0 \leq n < 10 \wedge P(n) \wedge Q(n)) \quad \text{y} \quad (\forall n: \mathbb{Z}) (0 \leq n < 10 \rightarrow Q(n))$
    \item $k=0 \wedge (\exists n:\mathbb{Z})(0 \leq n < 10 \wedge P(n) \wedge Q(n)) \quad \text{y} \quad k=0 \wedge (\forall n:\mathbb{Z})(0 \leq n < k \rightarrow Q(n))$
\end{enumerate}

\vspace{0.2cm}
{\color{gray}\hrule}
\vspace{0.4cm}

\textbf{Resolución:}

Recordemos del Ejercicio 5 que una proposición $\alpha$ es más fuerte que $\beta$ si la implicación $\alpha \rightarrow \beta$ es una tautología (siempre verdadera). Dicho de otra forma, siempre que $\alpha$ se cumple, $\beta$ está obligada a cumplirse.

\begin{enumerate}[label=\textbf{\alph*)}]

    \item Llamemos $\alpha = P(3)$ y $\beta = (\forall n: \mathbb{Z}) (0 \leq n < 5 \rightarrow P(n))$. \\
    - ¿$\alpha \rightarrow \beta$? Que $P$ valga para el número 3 no garantiza que valga para todos los números del 0 al 4. (No es tautología). \\
    - ¿$\beta \rightarrow \alpha$? Si sabemos que $P$ se cumple para \textbf{todos} los enteros entre 0 y 4, entonces obligatoriamente se cumple para el 3 (porque $0 \leq 3 < 5$). (Es tautología). \\
    \underline{\textbf{$\beta$ es más fuerte que $\alpha$.}}

    \vspace{0.4cm}
    \item Llamemos $\alpha = P(3)$ y $\beta = (\exists n: \mathbb{Z}) (0 \leq n < 5 \wedge P(n))$. \\
    - ¿$\alpha \rightarrow \beta$? Si sabemos que $P(3)$ es verdadero, entonces automáticamente \textbf{existe} al menos un número entre el 0 y el 4 que cumple $P$ (justamente, el 3). (Es tautología). \\
    - ¿$\beta \rightarrow \alpha$? Saber que \textit{algún} número en ese rango cumple $P$ no nos garantiza que sea específicamente el 3 (podría ser el 1, por ejemplo). (No es tautología). \\
    \underline{\textbf{$\alpha$ es más fuerte que $\beta$.}}

    \vspace{0.4cm}
    \item $\alpha$ afirma que \textit{si} un número en el rango cumple $P$, entonces cumple $Q$. $\beta$ afirma que \textbf{todos} los números en el rango cumplen $Q$ directamente (sin condiciones). \\
    - ¿$\alpha \rightarrow \beta$? Si solo sabemos que los que cumplen $P$ cumplen $Q$, podría haber números que no cumplan $P$ y tampoco $Q$. Así que $\beta$ no está obligada a ser cierta. \\
    - ¿$\beta \rightarrow \alpha$? Si sabemos como hecho absoluto que \textbf{todos} en el rango cumplen $Q$ ($\beta$ es verdadera), entonces cualquier condición que le pongamos adelante (como exigir $P(n)$) resultará en una implicación verdadera, porque el consecuente $Q(n)$ ya está garantizado. \\
    \underline{\textbf{$\beta$ es más fuerte que $\alpha$.}}

    \vspace{0.4cm}
    \item $\alpha$ afirma que existe al menos un número que cumple ambas ($P$ y $Q$). $\beta$ afirma que todos los números cumplen $Q$. \\
    - ¿$\alpha \rightarrow \beta$? Que un solo elemento cumpla todo no obliga a que \textbf{todos} los demás cumplan $Q$. \\
    - ¿$\beta \rightarrow \alpha$? Que todos cumplan $Q$ no garantiza que alguno cumpla $P$. Podría ocurrir que $P$ sea falso para todos, haciendo que $\alpha$ sea falsa. \\
    \underline{\textbf{Son incomparables (ninguna es más fuerte).}}

    \vspace{0.4cm}
    \item Este es el inciso trampa. Analicemos $\beta$: la segunda parte exige $(\forall n:\mathbb{Z})(0 \leq n < k \rightarrow Q(n))$. Como la primera parte fija $k=0$, el rango queda $0 \leq n < 0$. Este antecedente es siempre Falso, por lo que esa implicación universal es \textbf{siempre verdadera} por vacuidad. Entonces $\beta$ equivale lógicamente a $(k=0 \wedge True)$, que es simplemente $k=0$. \\
    La fórmula $\alpha$ equivale a $(k=0 \wedge \text{Existencial})$. \\
    - ¿$\alpha \rightarrow \beta$? Si ocurre $(k=0 \wedge \text{Existencial})$, es obvio que ocurre $k=0$. Es decir, $(X \wedge Y) \rightarrow X$ es una tautología. \\
    - ¿$\beta \rightarrow \alpha$? Si sabemos que $k=0$, no tenemos idea de si el existencial es verdadero o falso. \\
    \underline{\textbf{$\alpha$ es más fuerte que $\beta$.}}

\end{enumerate}

\vspace{0.5cm}
\hrule
\vspace{0.5cm}

\end{document}
